\documentclass[11pt]{amsart}
\usepackage{amsmath,amssymb,amsthm}
\usepackage[a4paper,margin=30mm]{geometry}
\usepackage{microtype}
\usepackage{enumitem}

\newtheorem{theorem}{Theorem}
\newtheorem{lemma}[theorem]{Lemma}
\newtheorem{corollary}[theorem]{Corollary}
\theoremstyle{definition}
\newtheorem{definition}[theorem]{Definition}
\theoremstyle{remark}
\newtheorem*{remark}{Remark}

\newcommand{\Q}{\mathbb{Q}}
\newcommand{\R}{\mathbb{R}}
\newcommand{\C}{\mathbb{C}}
\newcommand{\ran}{r_{\mathrm{an}}}
\newcommand{\Rad}{\mathcal{R}_{E,L}}
\newcommand{\Hs}{\mathcal{H}_L}
\DeclareMathOperator{\supp}{supp}

\title[Optimal explicit-formula bounds for analytic ranks]{Optimal explicit-formula bounds for analytic ranks:\\ the limiting radix kernel}
\author{Sabrina Hackfort (MEG)}
\address{Independent researcher (MEGALODON Research), Germany}
\email{meg.depth@proton.me}
\date{25 October 2026}



\begin{document}
\begin{abstract}
For an elliptic curve $E/\mathbb{Q}$ and a window $[-L,L]$, the radix kernel of Hackfort (2026) is the test function that
is optimal for each curve, within a finite-dimensional space, for bounding the analytic rank of $E$ under GRH by the
explicit formula. Assuming GRH, we prove that the optimum over all admissible test functions with this support exists, is
unique, equals the reciprocal of the value at the origin of the reproducing kernel of the Hilbert space defined by the Weil
quadratic form, and is the limit of the finite-dimensional radix bounds, with an exact error identity. By Akhiezer's
factorisation theorem, no continuous test function $g$ with $\supp g\subset[-2L,2L]$, $\hat g\ge0$ and $\hat g(0)>0$ gives a
smaller bound. The key step is a coercivity estimate for the Weil form, obtained by writing it as a truncated Wiener--Hopf
form with a logarithmically growing symbol.
\end{abstract}
\maketitle

\section*{Introduction}
Under GRH, the explicit formula bounds the analytic rank $\ran$ of $E$ by $W_E(g)/\hat g(0)$ for every admissible test
function $g$; Mestre \cite{Mestre} introduced the method and Bober \cite{Bober} made it practical with a fixed Fej\'er-type
kernel. In \cite{Hac} the test function is optimised for each curve within the span of $P_{2j}(x/L)$, $j<K$; the resulting
\emph{radix bound} $B_K(L)=1/(v^\top G^{-1}v)$, where $G$ is the Gram matrix of the Weil quadratic form, saturates
numerically already for $K\approx10$. Assuming GRH, this note proves that $B_K(L)$ converges to the best bound obtainable
from continuous test functions supported in $[-2L,2L]$ with non-negative Fourier transform, and identifies the limit and
the extremal function, the \emph{limiting radix kernel} $k_{E,L}$ (Theorem~\ref{thm:main}).


\noindent Throughout, $E/\Q$ is an elliptic curve of conductor $N$ and we assume GRH for $L(E,s)$, normalised with
centre $s=\tfrac12$; its zeros are $\tfrac12+i\gamma$. For $p\nmid N$ write $a_p/\sqrt p=\alpha_p+\beta_p$ with
$\alpha_p\beta_p=1$, $|\alpha_p|=1$, and put $c_{p^k}=\alpha_p^k+\beta_p^k$; for $p\mid N$ put $c_{p^k}=(a_p/\sqrt p)^k$.
Let $w_n=\Lambda(n)c_n/\sqrt n$, $\hat f(t)=\int f(x)e^{itx}\,dx$, $f^\star(x)=\overline{f(-x)}$.

\section{The Weil form as a Toeplitz form}

For $f\in L^2(\R)$ with support in $[-L,L]$ and $g=f*f^\star$ (so $\hat g=|\hat f|^2$, $\supp g\subset[-2L,2L]$),
the right-hand side of the explicit formula \cite[Sect.~2]{Hac} is
\[
  W_E(g)=g(0)\bigl(\log N-2\log2\pi\bigr)+\mathrm{Arch}(g)-2\sum_{n<e^{2L}}w_n\operatorname{Re}g(\log n),
\]
\[
  \mathrm{Arch}(g)=2\Bigl[-\gamma_{\mathrm E}\,g(0)+\int_0^{2L}\frac{g(0)-\operatorname{Re}g(u)}{e^u-1}\,du-g(0)\log\bigl(1-e^{-2L}\bigr)\Bigr],
\]
with values in $(-\infty,+\infty]$. The same formula defines $W_E(g)$ for every continuous $g$ with
$\supp g\subset[-2L,2L]$ and $\hat g\ge0$; then $g(0)\ge|g(u)|$, so the integrand in $\mathrm{Arch}(g)$ is non-negative.
For smooth $g$ the explicit formula gives $W_E(g)=\sum_\gamma\hat g(\gamma)$ \cite{Mestre,Bober}.

\begin{lemma}\label{lem:symbol}
For every such $f$,
\[
  W_E(g)=\frac1{2\pi}\int_\R|\hat f(t)|^2M(t)\,dt,
\]
\[
  M(t)=\log N-2\log2\pi+2\operatorname{Re}\psi(1+it)-2\sum_{n<e^{2L}}w_n\cos(t\log n),
\]
where $\psi=\Gamma'/\Gamma$. $M$ is real and even, and $|M(t)-2\log(e+|t|)|\le C_{E,L}$ for all $t\in\R$.
\end{lemma}

\begin{proof}
From $\psi(z)=-\gamma_{\mathrm E}+\int_0^\infty\frac{e^{-u}-e^{-zu}}{1-e^{-u}}\,du$ we get
$\operatorname{Re}\psi(1+it)=-\gamma_{\mathrm E}+\int_0^\infty\frac{1-\cos tu}{e^u-1}\,du$. Since
$g(u)=\frac1{2\pi}\int|\hat f|^2e^{-itu}dt$, we have $g(0)-\operatorname{Re}g(u)=\frac1{2\pi}\int|\hat f|^2(1-\cos tu)\,dt$, and
$g(u)=0$ for $u>2L$; Tonelli (the integrand is non-negative) gives
$\frac1{2\pi}\int|\hat f|^2\,2\operatorname{Re}\psi(1+it)\,dt=\mathrm{Arch}(g)$, including the case where both sides are $+\infty$.
The prime sum is finite. The bound follows from $\operatorname{Re}\psi(1+it)=\log|t|+O(t^{-2})$ for $|t|\ge1$.
\end{proof}

So $W_E$ is a truncated Wiener--Hopf (continuous Toeplitz) form with symbol $M$ on the Paley--Wiener space. The symbol is not non-negative, but under GRH
the form is. Let $\mu(t)=2\log(e+|t|)$ and let $\Hs$ be the space of $f\in L^2(\R)$ with support in $[-L,L]$ and
\[
  \|f\|_\mu^2=\frac1{2\pi}\int|\hat f|^2\mu<\infty,\qquad Q(f,h)=\frac1{2\pi}\int\hat f\,\overline{\hat h}\,M,\quad Q(f)=Q(f,f),
\]
and, once $Q$ is known to be positive definite, $\|f\|_Q=Q(f)^{1/2}$.
$\Hs$ is a Hilbert space, $|Q(f)|\le C\|f\|_\mu^2$, and $W_E(f*f^\star)=+\infty$ for $f\notin\Hs$. The indicator
$\mathbf 1_{[-L,L]}$ (which generates Bober's Fej\'er kernel $g=\mathbf 1_{[-L,L]}*\mathbf 1_{[-L,L]}$) and every polynomial restricted to $[-L,L]$ lie in $\Hs$.

\begin{lemma}\label{lem:compact}
The embedding $\Hs\hookrightarrow L^2(\R)$ is compact.
\end{lemma}
\begin{proof}
For $\|f\|_\mu\le1$ the functions $\hat f$ are uniformly bounded and uniformly Lipschitz ($|\hat f'|\le\int|xf|\le L\sqrt L$),
so a subsequence converges locally uniformly (Arzel\`a--Ascoli), while
$\frac1{2\pi}\int_{|t|>T}|\hat f|^2\le1/(2\log(e+T))$ uniformly; by Plancherel the subsequence is Cauchy in $L^2$.
\end{proof}

\begin{lemma}\label{lem:formula}
For $f\in\Hs$, $\sum_\gamma|\hat f(\gamma)|^2\le Q(f)$ (zeros with multiplicity); in particular $\ran|\hat f(0)|^2\le Q(f)$.
If $Q(f)=0$ then $f=0$.
\end{lemma}
\begin{proof}
For smooth $f$ this is the explicit formula (with equality) and Lemma~\ref{lem:symbol}. For general $f\in\Hs$, dilate
($f(x/(1-\delta))$) and mollify; the approximants are smooth with support in $[-L,L]$ and converge to $f$ in $\Hs$ (on the Fourier side the
dilations are uniformly bounded on $L^2(\R,\mu\,dt)$ since $\mu(s/(1-\delta))\le C\mu(s)$, and converge strongly on
$C_c(\R)$, which is dense there), hence
$Q$ converges, and $\hat f$ converges uniformly on $\R$. Fatou's lemma gives the inequality.

If $Q(f)=0$, then $\hat f(\gamma)=0$ for every ordinate. Under GRH the multiplicity of a zero at $\tfrac12+i\gamma$ is
$\ll\log(N(|\gamma|+3))/\log\log(N(|\gamma|+3))$: apply \cite[Thm.~5]{CCM} to the twist $L(s+i\gamma,E)$, which satisfies
GRH, the Ramanujan bound and a functional equation, with analytic conductor $\asymp N(|\gamma|+3)^2$. Since the number of
ordinates in $[-T,T]$ counted with multiplicity is $\frac T\pi\log\frac{NT^2}{(2\pi e)^2}+O(\log NT)$, there are
$\gg T\log\log T$ distinct ones. An entire function of exponential type has $O(R)$ zeros in $|z|\le R$ (Jensen), so
$\hat f\equiv0$.
\end{proof}

\begin{theorem}[coercivity]\label{thm:coercive}
There is $\lambda_{E,L}>0$ with $Q(f)\ge\lambda_{E,L}\|f\|_\mu^2$ for all $f\in\Hs$.
\end{theorem}
\begin{proof}
Write $M=\mu+R$ with $R$ bounded. Then $Q(f)=\langle(I+T)f,f\rangle_\mu$, where $\langle Tf,h\rangle_\mu=\frac1{2\pi}\int\hat f\,
\overline{\hat h}\,R$; $T$ factors through the compact embedding of Lemma~\ref{lem:compact}, so it is compact and
self-adjoint. By Riesz--Schauder \cite[Thm.~VI.15]{RS}, $\sigma(I+T)\cap(-\infty,\tfrac12)$ consists of finitely many
eigenvalues, each positive by Lemma~\ref{lem:formula}. Hence $\min\sigma(I+T)>0$.
\end{proof}

\section{The limiting radix kernel is optimal}

\begin{theorem}\label{thm:main}
Assume GRH for $L(E,s)$ and let $L>0$.
\begin{enumerate}[label=\textup{(\roman*)}]
\item There is a unique $k=k_{E,L}\in\Hs$ with $Q(f,k)=\hat f(0)$ for all $f\in\Hs$; it is real and even. Writing
      $\Rad(w,z)$ for the reproducing kernel of the space $\{\hat f:f\in\Hs\}$ with norm $\|\cdot\|_Q$, we have
      $\hat k(z)=\Rad(z,0)$ and $\Rad(0,0)=\hat k(0)=Q(k)>0$.
\item $\displaystyle B^*(L):=\inf_{0\ne f\in\Hs}\frac{Q(f)}{|\hat f(0)|^2}=\frac1{\Rad(0,0)}$, attained exactly by the
      multiples of $k$; in particular $\ran\le\lfloor B^*(L)\rfloor$.
\item For every continuous $g$ with $\supp g\subset[-2L,2L]$, $\hat g\ge0$ on $\R$ and $\hat g(0)>0$,
      \[
        \frac{W_E(g)}{\hat g(0)}\ \ge\ B^*(L).
      \]
      Thus among all explicit-formula bounds whose test function lies in this class, so that only the coefficients
      $c_n$ with $n<e^{2L}$ enter, $B^*(L)$ is the smallest, and the radix bounds $B_K(L)$ converge to it by (iv).
\item Let $B_K(L)=1/(v^\top G^{-1}v)$ be the bound computed in the basis $P_{2j}(x/L)$, $j<K$ \cite[Sect.~3]{Hac}. Then $B_K(L)$
      is non-increasing in $K$, $B_K(L)\to B^*(L)$, and
      \[
        B_K(L)-B^*(L)=\frac{\|k-P_Kk\|_Q^2}{\Rad(0,0)\,\|P_Kk\|_Q^2},
      \]
      where $P_K$ is the $Q$-orthogonal projection onto $V_K=\operatorname{span}\{P_{2j}(x/L):j<K\}$.
\end{enumerate}
\end{theorem}

\begin{proof}
(i) By Theorem~\ref{thm:coercive}, $(\Hs,Q)$ is a Hilbert space, and $\ell(f)=\hat f(0)=\langle f,\mathbf 1_{[-L,L]}\rangle$
is continuous, with $\langle\cdot,\cdot\rangle$ the inner product of $L^2(\R)$ ($|\ell(f)|\le\sqrt L\|f\|_\mu$). Riesz representation gives $k$. Since $M$ is even and real, $k(-x)$ and
$\overline{k}$ represent $\ell$ as well, so $k$ is real and even. $Q(k)=\ell(k)$, and $Q(k)>0$ as $\ell\ne0$.

(ii) Cauchy--Schwarz in $(\Hs,Q)$: $|\hat f(0)|^2=|Q(f,k)|^2\le Q(f)Q(k)$, with equality iff $f\in\C k$. The bound on
$\ran$ is Lemma~\ref{lem:formula}.

(iii) As $g$ is continuous at $0$ and $\hat g\ge0$, $\hat g\in L^1(\R)$ (Fej\'er summation at $0$). $\hat g$ is entire of exponential type $\le2L$,
bounded and non-negative on $\R$. Hence $\int\log^+\hat g(t)\,(1+t^2)^{-1}dt<\infty$, and by Akhiezer's theorem \cite{Akhiezer,Boas}, $\hat g=|F|^2$ with $F$ entire of exponential
type $\le L$; $F\in L^2(\R)$, so $F=\hat f$ with $\supp f\subset[-L,L]$ (Paley--Wiener), and $g=f*f^\star$. If
$f\notin\Hs$, then $W_E(g)=+\infty$. Otherwise $W_E(g)=Q(f)$ by Lemma~\ref{lem:symbol}, $\hat g(0)=|\hat f(0)|^2$, and (ii)
applies.

(iv) $G$ is the Gram matrix of $Q$ on $V_K$ (invertible since $Q$ is positive definite) and $v_j=\ell(P_{2j}(\cdot/L))$, so $v=(2L,0,\dots,0)$ by orthogonality of the Legendre polynomials. On $V_K$ the functional $\ell$ is represented by $P_Kk$, so $B_K=1/\|P_Kk\|_Q^2$ as in (ii); Pythagoras gives
the identity, and $B_K$ is non-increasing because the $V_K$ are nested. Even polynomials are dense in $\Hs\cap\{f\text{
even}\}$: for $h\in C^1[-L,L]$, integration by parts gives $|\widehat{h\mathbf 1_{[-L,L]}}(t)|\ll\|h\|_{C^1}/|t|$, hence
$\|h\mathbf 1_{[-L,L]}\|_\mu\ll\|h\|_{C^1}$; smooth even functions supported in $(-L,L)$ are dense, and are $C^1$-limits of
even polynomials. With the norm equivalence of Theorem~\ref{thm:coercive}, $\|k-P_Kk\|_Q\to0$.
\end{proof}

\begin{remark}
Part (iii) concerns certificates that use $\hat g\ge0$ on all of $\R$, which is what one needs when the zeros are unknown.
Certificates using further information (for example the parity or integrality of $\ran$, or signs of individual prime
terms) are not covered. In \cite{Hac}, on the curve \#4239 of the ICARM Elliptic Curve Rank Leaderboard \cite{ICARM} (rank $\ge32$), increasing $K$
from $10$ to $20$ changes the bound only in the fifth decimal; by (iv) this is consistent with $\|k-P_Kk\|_Q$ being small
for $K\ge10$. A
quantitative rate in (iv) would follow from regularity of $k$, which we do not pursue here.
\end{remark}

\subsection*{Tools}
PARI/GP, Arb and Rust (the radix-rank code of \cite{Hac}) for the numerical values quoted from \cite{Hac}; Claude (Anthropic).
The author takes full responsibility for the content; all numerical results are reproducible from the code and
certificates of \cite{Hac}.

\begin{thebibliography}{9}
\bibitem{Akhiezer} N.\,I.~Akhiezer, \emph{On the theory of entire functions of finite degree}, Doklady Akad.\ Nauk SSSR (N.S.) \textbf{63} (1948), 475--478.
\bibitem{Boas} R.\,P.~Boas, \emph{Entire Functions}, Academic Press, 1954, pp.~124--132.
\bibitem{Bober} J.\,W.~Bober, \emph{Conditionally bounding analytic ranks of elliptic curves}, ANTS X, Open Book Ser.\ \textbf{1} (2013), 135--144.
\bibitem{CCM} E.~Carneiro, V.~Chandee, M.\,B.~Milinovich, \emph{A note on the zeros of zeta and $L$-functions}, Math.\ Z.\ (2015), doi:10.1007/s00209-015-1485-9; arXiv:1503.00955.
\bibitem{ICARM} ICARM, \emph{Elliptic Curve Rank Leaderboard}, \texttt{https://elliptic-rank.icarm.cloud}, curve \#4239 (accessed 25 October 2026).
\bibitem{Hac} S.~Hackfort, \emph{Optimal windowed test functions for bounding analytic ranks of elliptic curves}, Hexagon Math, hexagon:2610.00346 (2026); code and certificates: \texttt{github.com/sectio-aurea-q/radix-rank}, doi:10.5281/zenodo.23299461.
\bibitem{Mestre} J.-F.~Mestre, \emph{Formules explicites et minorations de conducteurs de vari\'et\'es alg\'ebriques}, Compositio Math.\ \textbf{58} (1986), 209--232.
\bibitem{RS} M.~Reed, B.~Simon, \emph{Methods of Modern Mathematical Physics I: Functional Analysis}, Academic Press, 1980.
\end{thebibliography}
\end{document}
